\documentclass[11pt]{article}
\usepackage[a4paper,margin=26mm]{geometry}
\usepackage{amsmath,amssymb,amsthm,mathtools,hyperref}
\newcommand{\ddc}{dd^c}
\newcommand{\dc}{d^c}
\newcommand{\kap}{\kappa_\lambda}
\newtheorem{theorem}{Theorem}
\newtheorem{proposition}[theorem]{Proposition}
\newtheorem{lemma}[theorem]{Lemma}
\theoremstyle{remark}\newtheorem{remark}[theorem]{Remark}
\title{Lee curvature with constant minimum spectrum\newline
on compact locally conformally K\"ahler manifolds}
\author{}\date{}
\begin{document}\maketitle
\noindent\textbf{Working draft.} Complete independent review and
reconstruction are pending. The unrestricted P03 conjecture is not
resolved here.
\begin{abstract}
We prove that a compact strictly locally conformally K\"ahler metric
whose Lee curvature has constant least relative eigenvalue has
semipositive Lee curvature. The least eigenvalue is necessarily zero. The proof uses a strictly increasing volume polynomial
and compact Stokes, and applies in particular to compact quotients
with invariant LCK tensors. An explicit compact family has constant
least eigenvalue and nonconstant remaining spectrum, so the condition
is strictly weaker than constant spectrum. No classification is used. We give an explicit compact Hopf metric with indefinite
Lee curvature and a smooth correction to a semipositive representative.
In that example the first degeneracy of the volume ray leaves positive
total volume, explaining the limitation of the constant-spectrum
argument. All statements retain ordinary Bott--Chern cohomology.
\end{abstract}

\section{Conventions and the restricted statement}
Let $X$ be connected compact complex of dimension $n\geq2$. A smooth
positive real $(1,1)$ form $\omega$ is strictly LCK when
\[
 d\omega=\lambda\wedge\omega,\qquad d\lambda=0,
 \qquad [\lambda]\ne0\text{ in }H^1_{\rm dR}(X,\mathbb R).
\]
Use ordinary wedge powers and
\[
 \dc=\frac{i}{4\pi}(\bar\partial-\partial),\qquad
 \ddc=\frac{i}{2\pi}\partial\bar\partial,\qquad
 \lambda^c=\frac{i}{4\pi}(\lambda^{0,1}-\lambda^{1,0}),
 \qquad\kap=-d\lambda^c.
\]
Locally $\lambda=df$, so $\kap=-\ddc f$. Thus $\kap$ is smooth,
real, closed and of type $(1,1)$, and is de Rham exact. Its ordinary
Bott--Chern class need not vanish. Pseudoeffectivity means existence
of a closed positive current in that class.

\begin{theorem}\label{th:constant}
Suppose the LEAST eigenvalue of $\kap$ relative to $\omega$ is
constant on $X$. Then $\kap$ is semipositive, its least eigenvalue
equals zero, and its ordinary Bott--Chern class is nonzero and
pseudoeffective. No constancy of the other eigenvalues or of the
rank is assumed.

If the entire spectrum is constant, the rank is constant in
$\{1,\ldots,n-1\}$. In particular the conclusion applies to any
actual compact quotient $\Gamma\backslash G$ with left-invariant
complex structure and left-invariant strictly LCK tensors
$\omega,\lambda$.
\end{theorem}

The minimum-spectrum hypothesis is a condition on this metric and
its Lee form, rather than merely on the complex manifold. The
unrestricted assertion for every compact strictly LCK manifold is
the conjecture recorded in \cite[Conjecture 5.13]{OV}. Neither a
single positive Gauduchon pairing nor the restricted theorem proves
that assertion. No literature priority is claimed.

\section{Compact Stokes and the volume ray}
The two-form
\[
 Q=\lambda\wedge\lambda^c
   =\frac{i}{2\pi}\lambda^{1,0}\wedge\lambda^{0,1}
\]
is semipositive and is nonzero on an open set. The latter follows
because the nonexact smooth form $\lambda$ is not identically zero.

\begin{lemma}\label{lem:stokes}
For $1\leq k\leq n$ one has
\[
 \int_X\kap^k\wedge\omega^{n-k}
  =(n-k)\int_X Q\wedge\kap^{k-1}\wedge\omega^{n-k}.
\]
Consequently
\[
 \int_X\kap\wedge\omega^{n-1}>0,\qquad
 \int_X\kap^n=0.
\]
The ordinary Bott--Chern class of $\kap$ is nonzero.
\end{lemma}
\begin{proof}
Use $d\kap=0$, $d\lambda^c=-\kap$ and
$d\omega^j=j\lambda\wedge\omega^j$. Differentiation gives
\[
 d(\lambda^c\wedge\kap^{k-1}\wedge\omega^{n-k})
 =-\kap^k\wedge\omega^{n-k}
   +(n-k)Q\wedge\kap^{k-1}\wedge\omega^{n-k}.
\]
Integration and smooth Stokes prove the identity. For $k=1$, the
right side is strictly positive since $n\geq2$, $\omega$ is positive
and $Q$ is nonzero on an open set. For $k=n$ it is zero.

If $\kap=\ddc a$ for a smooth real global $a$, the conformal LCK
metric $e^a\omega$ has Lee form $\lambda+da$ and zero Lee curvature.
Its Lee form remains nonexact. Applying the positive $k=1$ identity
to this metric is a contradiction. This concerns ordinary
Bott--Chern exactness, rather than de Rham exactness.
\end{proof}

\begin{proposition}\label{prop:ray}
On the interval where $\eta_t=\omega+t\kap$ is positive, the
polynomial $V(t)=\int_X\eta_t^n$ has degree at most $n-1$ and
\[
 V'(t)=n(n-1)\int_X Q\wedge\omega\wedge\eta_t^{n-2}>0.
\]
No Gauduchon condition on $\omega$ or $\eta_t$ is assumed.
\end{proposition}
\begin{proof}
Since $d\eta_t=d\omega=\lambda\wedge\omega$, smooth Stokes gives
\[
 0=\int_X d(\lambda^c\wedge\eta_t^{n-1})
  =-\int_X\kap\wedge\eta_t^{n-1}
    +(n-1)\int_X Q\wedge\omega\wedge\eta_t^{n-2}.
\]
Differentiating the finite wedge polynomial proves the formula.
All factors in the final integral are positive or semipositive,
and it is strictly positive on the open set where $\lambda\ne0$.
Lemma~\ref{lem:stokes} makes its top coefficient
$\int_X\kap^n$ vanish.
\end{proof}

\begin{proof}[Proof of Theorem~\ref{th:constant}]
Let $r$ be the constant least relative eigenvalue. If $r<0$, put
$t_0=-1/r>0$. For $0\leq t<t_0$ the form
$\eta_t=\omega+t\kap$ is positive. At $t_0$ it is semipositive
and has a kernel at EVERY point, independently of the other
eigenvalues. Thus $V(t_0)=0$. Proposition~\ref{prop:ray} and
polynomial continuity instead give $V(t_0)\geq V(0)>0$.
This contradiction gives $r\geq0$ and pointwise semipositivity.

The continuous top power $\kap^n$ is now nonnegative and has zero
integral by Lemma~\ref{lem:stokes}. It vanishes identically. Thus
the least eigenvalue is zero at every point. The same lemma proves
that the ordinary Bott--Chern class is nonzero, while $\kap$ itself
is its closed positive representative.

If all eigenvalues are constant, the rank is constant. It cannot
be zero, and cannot be $n$, hence lies between one and $n-1$.
For invariant tensors on an actual compact quotient, the relative
matrix is constant in a left-invariant frame. No lattice existence
for an abstract Lie algebra is inferred from this calculation.
\end{proof}

\begin{remark}
Every invariant Hermitian metric $\gamma$ on such a compact quotient
is Gauduchon: $\ddc\gamma^{n-1}$ is an invariant top form, hence a
constant multiple of the invariant volume. Smooth Stokes makes
that constant zero. The theorem uses pointwise semipositivity of
$\kap$, so it controls pairings against all positive Gauduchon
tests, including those which are not invariant.
\end{remark}

\section{An indefinite Hopf representative and its correction}
Fix $0<\varepsilon<1$, and on $\mathbb C^n\setminus\{0\}$ put
\[
 \rho=\sum_{j=1}^n|z_j|^2,\qquad
 \Phi=\rho+\varepsilon\operatorname{Re}(z_1^2),\qquad
 \Omega=\ddc\Phi=\ddc\rho>0.
\]
The bounds $(1-\varepsilon)\rho\leq\Phi\leq(1+\varepsilon)\rho$
make $\Phi$ positive. The action $g(z)=z/2$ is free and properly
discontinuous. Polar coordinates show that its quotient $X$ is
connected compact, with a fundamental interval of logarithmic radii.
Both $\Phi$ and $\Omega$ scale by $1/4$ under $g$. Therefore
\[
 \omega=\Omega/\Phi,\qquad \lambda=-d\log\Phi
\]
descend and satisfy $d\omega=\lambda\wedge\omega$. The Lee period
along the quotient loop lifted from $z$ to $z/2$ is $2\log2$, so this
is a strictly LCK metric. Its Lee curvature is $\kap=\ddc\log\Phi$.

Writing $F_i=\bar z_i+\varepsilon z_1\delta_{i1}$, its complex
Hessian is
\[
 H_{i\bar j}=\frac{\delta_{ij}}\Phi-
                   \frac{F_i\overline{F_j}}{\Phi^2}.
\]
At $(1,0,\ldots,0)$ the diagonal entries are
$-\varepsilon/(1+\varepsilon)$ and $1/(1+\varepsilon)$ in the
remaining directions. Relative to $\omega$, its eigenvalues there
are $-\varepsilon,1,\ldots,1$. Thus this metric has indefinite Lee
curvature; a positive strictly psh automorphic potential can have
a logarithm which is not psh.

Nevertheless the smooth bounded deck-invariant function
\[
 u=\log\rho-\log\Phi
\]
descends and gives
\[
 \kap+\ddc u=\ddc\log\rho\geq0.
\]
Indeed the Hessian of $\log\rho$ is
$\delta_{ij}/\rho-\bar z_i z_j/\rho^2$, semipositive by
Cauchy--Schwarz. This is an actual smooth semipositive representative
of the ordinary Bott--Chern class. The example is not a
counterexample to pseudoeffectivity.

To examine its first volume-ray boundary, write $z_1=x+iy$ and
$R=\sum_{j>1}|z_j|^2$. The remaining relative curvature eigenvalue is
\[
 h=\frac{\varepsilon((1-\varepsilon)y^2-(1+\varepsilon)x^2)}
            {(1+\varepsilon)x^2+(1-\varepsilon)y^2+R},
 \qquad -\varepsilon\leq h\leq\varepsilon.
\]
The other $n-1$ eigenvalues are one. Hence the first positive-ray
boundary is $t_0=1/\varepsilon$. At this boundary
\[
 1+h/\varepsilon=
 \frac{2(1-\varepsilon)y^2+R}{\Phi}.
\]
It vanishes only on the real $z_1$ rays with $z_j=0$ for $j>1$.
The top power of $\eta_{t_0}$ is positive on a nonempty open set,
so its total volume is positive. First degeneracy on a general
compact LCK manifold therefore need not collapse the total volume.

\begin{remark}
The proof of the restricted theorem uses a globally constant least
eigenvalue to force degeneration everywhere at one parameter. The Hopf
calculation explains why that step does not prove the unrestricted
conjecture. No classification, all-cover construction, unbounded
potential assertion or literature priority is established here.
\end{remark}

\section{Constant minimum with genuinely variable remaining spectrum}
The weaker hypothesis is nonvacuous beyond constant spectrum.
On the SAME scalar Hopf quotient, put $s=\rho$, $q=|z_1|^2$,
and, for any $0<\varepsilon<1$, define
\[
 \Psi=s+\varepsilon\frac{q^2}{s},\qquad
 \Omega_\varepsilon=\ddc\Psi,\qquad
 \omega_\varepsilon=\Omega_\varepsilon/s,\qquad
 \lambda=-d\log s.
\]
Here $\Psi$ is smooth on the punctured space and scales by $1/4$
under $z\mapsto z/2$. Let $a=q/s\in[0,1]$, and use subscripts
for complex derivatives. Direct differentiation gives
\[
 \left(\frac{q^2}{s}\right)_{i\bar j}
 =\frac{2}{s}(q_i-a s_i)(q_{\bar j}-a s_{\bar j})
       +2a\,\delta_{i1}\delta_{j1}-a^2\delta_{ij}.
\]
The first two matrix terms are positive semidefinite, so the
complex Hessian is at least $-I$. Consequently
$\Omega_\varepsilon\geq(1-\varepsilon)\ddc s>0$.
This is an actual positive closed automorphic form, not merely
an algebraic assignment of eigenvalues. The form
$\omega_\varepsilon$ descends to the compact quotient, with
$d\omega_\varepsilon=\lambda\wedge\omega_\varepsilon$.
The Lee period remains $2\log2$, proving strict LCK.

Its Lee curvature is $\kap=\ddc\log s$. This is semipositive
of rank $n-1$, so the least relative eigenvalue is identically
zero for every member of this family. At $z=(1,0,\ldots,0)$
the complex Hessian of $q^2/s$ is
$\operatorname{diag}(1,-1,\ldots,-1)$, and the relative
curvature spectrum is
\[
 0,\ (1-\varepsilon)^{-1},\ldots,(1-\varepsilon)^{-1}.
\]
At $z=(0,1,0,\ldots,0)$, $q^2/s$ vanishes to fourth order,
so its complex Hessian is zero. The relative spectrum there
is $0,1,\ldots,1$. Thus the remaining spectrum is genuinely
nonconstant for every $\varepsilon>0$, in every $n\geq2$.
The constant-minimum condition is strictly weaker than requiring
all eigenvalues constant. This is a sufficiency example, not a
counterexample to the unrestricted conjecture.

\paragraph{Verification boundary.}
The original compact strictly LCK hypotheses and ordinary
Bott--Chern class are retained. The minimum-spectrum theorem
and both compact families above are new unverified author
statements requiring independent reconstruction. The earlier
constant-spectrum working draft and its frozen review inputs
are preserved; they are not correctness grants. No unrestricted
P03 closure or literature priority is asserted.

\begin{thebibliography}{9}
\bibitem{OV} L. Ornea and M. Verbitsky,
\emph{The Lee--Gauduchon cone on complex manifolds},
\href{https://arxiv.org/abs/2411.05595v2}{arXiv:2411.05595v2},
Conjecture 5.13 (problem context).
\end{thebibliography}
\end{document}
